
\subsection{Etapas sistemáticas}
Para abordar el desafío, se propone una secuencia de cinco etapas sistémicas y estructuradas que permiten comprender, delimitar y proyectar la problemática con el objetivo de desarrollar una solución adecuada frente a este sistema complejo.\\

\begin{itemize}
    \item \textbf{Etapa 1: Levantamiento y comprensión del contexto.}
    La primera etapa consiste en reconocer el sistema general en el que se inserta la problemática. Su propósito es construir una visión amplia del entorno hospitalario interno, identificando los actores del flujo asistencial, los canales de comunicación existentes entre turnos, las herramientas formales e informales que el personal utiliza para traspasar información y las restricciones normativas y operativas que condicionan cualquier intervención de diseño.

    \item \textbf{Etapa 2: Caracterización del usuario y comprensión de la experiencia actual.}
    La segunda etapa se enfoca en profundizar en el usuario principal del sistema, que es el profesional asistencial que registra, consulta y traspasa el contexto clínico del paciente. El objetivo es comprender su experiencia real en el cambio de turno: qué información busca, cómo la obtiene hoy, qué le falta, qué le genera fricción y bajo qué condiciones físicas y cognitivas opera.

    \item \textbf{Etapa 3: Definición del estado de conflicto y formulación del problema de diseño.}
    Esta etapa tiene como objetivo modelar la situación actual que atraviesa el usuario, identificando las tensiones del sistema para formular de manera precisa el problema de diseño. A partir de esta definición, se busca delimitar con claridad la dimensión específica del problema que será abordada por el proyecto, diferenciándola de otros problemas relacionados pero fuera del alcance de esta intervención.

    \item \textbf{Etapa 4: Determinación de criterios.}
    En esta etapa se establecen los criterios que orientarán la propuesta, considerando aspectos de registro ágil, trazabilidad, continuidad asistencial, no interferencia con la labor clínica, privacidad del dato y usabilidad en condiciones de alta carga cognitiva. Su propósito es traducir los hallazgos previos en lineamientos concretos que permitan orientar el desarrollo de una solución coherente con el contexto y con las necesidades detectadas.

    \item \textbf{Etapa 5: Proyección conceptual.}
    La última etapa consiste en formular soluciones parciales y conceptos iniciales que respondan al sistema de necesidades detectado. Esta etapa permite proyectar direcciones de solución a partir de la comprensión obtenida en las fases anteriores, estableciendo una base conceptual para el desarrollo posterior del sistema y de sus interfaces.
\end{itemize}
